\begin{table}[H]
\centering
\fontsize{8.5}{10.5}\selectfont
\renewcommand{\arraystretch}{1.15}
\setlength{\tabcolsep}{4pt}
\caption{Síntesis comparativa por perfil y consulta. Las razones usan la mediana; no se declara un ganador cuando los IQR se solapan.}
\label{tab:sintesis_colas}
\begin{tabular}{@{}lrrrrr@{}}
\toprule
\textbf{Perfil} & \textbf{4-ary/Binary} & \textbf{Fibonacci/Binary} & \textbf{Gana 4-ary} & \textbf{Gana Binary} & \textbf{IQR solap.} \\
\midrule
Disperso ($|E|=4|V|$) & 0.964 & 1.345 & 3 de 5 & 2 de 5 & 5 de 5 \\
Denso ($|E|=0.08|V|(|V|-1)$) & 0.954 & 1.374 & 4 de 5 & 1 de 5 & 5 de 5 \\
\midrule
Total sintético & 0.959 & 1.360 & 7 de 10 & 3 de 10 & 10 de 10 \\
Consulta OSM & 1.269 & 1.666 & 0 & 1 & -- \\
\bottomrule
\end{tabular}
\par\vspace{0.1cm}
\raggedright\fontsize{8}{10}\selectfont\textit{Nota.} Una razón 4-ary/Binary menor que 1 indica que el 4-ario fue más rápido. La columna IQR se refiere a 4-ary frente a Binary. Fibonacci no ganó ninguna escala sintética.
\end{table}
